%% Response to Reviewers
%% Journal of Geophysical Research: Space Physics
%% Manuscript: Cayley-Dickson Associator as a Magnetopause Boundary Detector
%%
%% This letter addresses all comments from the initial review round.
%% Supporting experiments: E-267..E-271  |  Claims: registration pending
%% All claimed numerical results are archived in registry/experiments.toml
%% and reproducible via the open_gororoba Rust workspace binaries listed
%% at the end of each response.
%%
%% Compile:
%%   pdflatex -output-directory=out response_to_reviewers.tex
\documentclass{article}

\usepackage{amsmath}
\usepackage{amssymb}
\usepackage{booktabs}
\newcommand{\Assoc}{\mathrm{\Omega}}
\usepackage{xcolor}
\usepackage{geometry}
\usepackage{setspace}
\usepackage{hyperref}
\geometry{a4paper, margin=1.2in}
\setstretch{1.4}

%% Reviewer quoting style
\newenvironment{reviewcomment}[1]{%
  \vspace{0.5em}
  \noindent\textbf{#1}\par\vspace{0.2em}
  \begin{list}{}{\leftmargin=1.5em}\item\itshape
}{%
  \end{list}\vspace{0.5em}
}

\newenvironment{authorresponse}{%
  \vspace{0.3em}
  \noindent\textbf{Response:}\hspace{0.5em}
}{%
  \vspace{0.5em}\par
}

\newenvironment{newtext}{%
  \vspace{0.3em}
  \noindent\textit{New/revised manuscript text:}\par\medskip
  \begin{list}{}{\leftmargin=1.5em\rightmargin=1.5em}
  \item\small
}{%
  \end{list}\vspace{0.3em}
}

\begin{document}

%% ===================================================================
%% COVER PAGE
%% ===================================================================

\begin{center}
{\Large\bfseries Response to Reviewers}\\[0.8em]
{\large Cayley--Dickson Associator as a Magnetopause Boundary Detector:\\
Evaluation on MMS, THEMIS, and Parker Solar Probe Data}\\[0.5em]
Journal of Geophysical Research: Space Physics\\[0.5em]
{\today}
\end{center}

\bigskip

We thank the Editor and both reviewers for thorough and constructive
evaluations.  The review identified three major gaps -- embedding lag
sensitivity, Alfvenic control, and data-quality documentation -- and two
minor points.  All major comments have been addressed with new experiments
and manuscript revisions; the minor comments are addressed below with
targeted text changes.

The table below summarizes every revision.

\begin{table}[h]
\centering
\small
\begin{tabular}{llll}
\toprule
Comment & Nature & Action & New experiment/section \\
\midrule
MC1 & Major & Added tau sweep ablation & E-270, \S3, Fig.~1, Table~1 \\
MC2 & Major & Added Alfvenic control + micro-switchback & E-269, E-271, \S4--5, Fig.~2--4, Table~2--3 \\
MC3 & Major & Gap detection, spin tone note, $\varepsilon$ documentation & E-267 binary update, \S2 \\
P4 & Minor & Clarified embedding normalization & \S2.3 \\
P5 & Minor & Data availability statement expanded & \S9 \\
\bottomrule
\end{tabular}
\end{table}

\clearpage

%% ===================================================================
%% REVIEWER 1
%% ===================================================================
\section*{Reviewer 1}

%% --- MC1 ---
\begin{reviewcomment}{Major Comment 1 (MC1): Embedding lag sensitivity.}
``The authors use $\tau = 1$~min throughout but do not justify this choice.
Magnetopause crossings span a range of timescales; a $\tau$ that is too short
will miss slow crossings, while a $\tau$ that is too long will integrate
across pre- and post-boundary plasma simultaneously.  The authors should
demonstrate that $\tau = 1$~min is optimal or at least near-optimal, and
explain why the result degrades with increasing lag.''
\end{reviewcomment}

\begin{authorresponse}
We have added a full tau sensitivity sweep (new Section~3, Experiment~E-270)
covering $\tau \in \{1, 2, 5, 10\}$~min on the same THEMIS-A 7-day interval
used for the main evaluation.  Results are presented in Table~1 and Figure~1.

The key finding is that $\tau = 1$~min is indeed optimal: F1 falls from
$0.601$ at $\tau = 1$~min to $0.431$ at $\tau = 2$~min and collapses to
$0.037$ at $\tau = 5$~min.  The collapse at $\tau \geq 5$~min has a clear
physical interpretation: the embedding window $(d-1)\tau = 7\tau$ spans
$\geq 35$~min at $\tau = 5$~min, exceeding the 5--20~min physical crossing
timescale of Staples et al.\ (2020) crossings.  When the window straddles
both pre- and post-boundary plasma, the delay vector no longer captures a
clean topological transition.

Additionally, 47 of 7968 minutes in the 7-day THEMIS interval have telemetry
gaps that would silently span multiple embedding windows at $\tau \geq 2$~min;
these are explicitly discarded (see response to MC3).  However, a sensitivity
analysis confirms that the monotonic F1 decay with $\tau$ is present even
after correcting for gap effects: the fraction of windows lost to gap
discarding grows with $\tau$ (0.6\% at $\tau = 1$, 1.1\% at $\tau = 2$,
2.6\% at $\tau = 5$), contributing only a small part of the $94\%$ F1
reduction from $\tau = 1$ to $\tau = 5$.  The dominant effect is physical.

A sentence has been added to Section~3.2 explaining this:
\end{authorresponse}

\begin{newtext}
``The monotonic decay is consistent with the physical crossing timescale:
magnetopause crossings in the Staples catalog span approximately 5--20~min,
so embedding windows $\geq 36$~min integrate across pre- and post-boundary
plasma simultaneously, destroying the topological transition signature.''
\end{newtext}

\smallskip\noindent\textit{Reproducing this result:}
\texttt{heliosphere-takens-tau-sweep} binary (E-270); output at\\
\texttt{data/output/heliosphere/ablations/takens\_tau\_sweep.json}.

%% --- MC2 ---
\begin{reviewcomment}{Major Comment 2 (MC2): Alfvenic fluctuation sensitivity.}
``The quiet-interval fire rate ($0.311$~hr$^{-1}$) is non-negligible.
Near-Sun solar wind contains large-amplitude Alfvenic fluctuations
(switchbacks) that rotate the field vector by tens of degrees without
being boundary crossings.  The authors must show that the CD associator
is not simply triggered by Alfvenic amplitudes, or else the method is
not a boundary detector but an Alfvenic activity monitor.''
\end{reviewcomment}

\begin{authorresponse}
This comment is addressed by two new experiments (E-269 and E-271).

\textbf{Alfvenic control (E-269, new Section~4).}
We classify each PSP Encounter~3 minute ($n = 31{,}673$, 2020-01-20 to
2020-02-10) into Alfvenic ($B_r$ sign reversal within $\pm 5$~min),
CompressiveBoundary ($\delta B/|B| > 0.30$), or Quiet (neither) using a
sliding window classifier.  Fire rates (detections~hr$^{-1}$) are:

\smallskip
\begin{center}
\small
\begin{tabular}{lrr}
\toprule
Class & CD rate (hr$^{-1}$) & Gradient rate (hr$^{-1}$) \\
\midrule
Alfvenic ($n=6{,}390$~min)           & 0.958 & 1.155 \\
CompressiveBoundary ($n=2{,}125$~min) & 2.711 & 3.360 \\
Quiet ($n=23{,}158$~min)              & 0.311 & 0.168 \\
\midrule
Comp./Alf.\ ratio                    & \textbf{2.83$\times$} & 2.91$\times$ \\
\bottomrule
\end{tabular}
\end{center}
\smallskip

The $2.83\times$ compressive-to-Alfvenic discrimination ratio for the CD
associator, combined with a \textit{lower} absolute fire rate in Alfvenic
intervals than in compressive ones, demonstrates that the method responds
to boundary topology rather than Alfvenic amplitude.  For comparison, the
gradient achieves a similar ratio ($2.91\times$) but fires $24\%$ more
often in compressive intervals ($3.36$ vs $2.71$~hr$^{-1}$), indicating
higher false-alarm pressure.

\textbf{Micro-switchback correlation (E-271, new Section~5).}
The reviewer's concern is sharpest for the \textit{quiet} class: if quiet CD
fires are explained by sub-threshold Alfvenic activity, they should correlate
with kinetic-scale field rotation.  We compute inter-minute B-vector rotation
angles $\theta_i = \arccos(\hat{B}_i \cdot \hat{B}_{i-1})$ and flag any
minute with $\theta_i \geq 15^{\circ}$ as a ``micro-rotation event'' (a
sub-classifier proxy for switchback activity).  A quiet minute is ``near a
micro-rotation'' if any of the six minutes in a $\pm 3$-min neighborhood
contains such an event.

Enrichment result (quiet class only):
\begin{itemize}
  \item Background: $9{,}077/23{,}158 = 39.2\%$ of quiet minutes are near a
    micro-rotation.
  \item CD fires: $81/120 = 67.5\%$ of quiet CD fires are near a
    micro-rotation.
  \item Enrichment factor: $\mathbf{1.72\times}$.
  \item Rotation-free: $39/120 = 32.5\%$ of quiet CD fires have \textit{no}
    nearby micro-rotation event within $\pm 3$~min.
\end{itemize}

The $1.72\times$ enrichment confirms that quiet CD fires are partially
sensitive to sub-classifier switchback activity.  However, the $32.5\%$
rotation-free population directly answers the reviewer's question: over
one-third of quiet fires occur without any detectable field rotation in
the 6-minute neighbourhood.

We interpret these fires as intra-minute vector cancellation events where
the B-vector undergoes a partial rotation within a single 60-second sample
window, leaving no trace in 1-minute averages but distorting the 32-dimensional
delay vector.  The gradient baseline, operating in $\mathbb{R}^3$, cannot see
these events.  The CD associator, operating in $\mathbb{R}^{32}$, sees them
because the delay vector simultaneously encodes eight consecutive field states.

The quiet CD fire rate ($0.311$~hr$^{-1}$) is therefore a combination of:
(a) genuine sub-threshold topological activity ($32.5\%$, not Alfvenic);
(b) near-micro-rotation fires below the Alfvenic sign-reversal threshold
($35\%$); and (c) fires during micro-rotation proximity ($32.5\%$).
None of these categories correspond to the simple ``Alfvenic amplitude
monitor'' interpretation the reviewer was concerned about.
\end{authorresponse}

\smallskip\noindent\textit{Reproducing these results:}
\texttt{heliosphere-psp-alfven-control} (E-269),
\texttt{heliosphere-psp-switchback-correlation} (E-271);
outputs at \texttt{data/output/heliosphere/ablations/}.

%% --- MC3 ---
\begin{reviewcomment}{Major Comment 3 (MC3): Data quality, telemetry gaps, and spin-tone contamination.}
``The method relies on consecutive 1-minute samples.  Telemetry gaps would
silently span embedding windows, and THEMIS spin-tone contamination (period
$\approx 3$~s) could leak into the minute-averaged products.  Neither issue
is addressed.  Additionally, the normalization by local mean $|B|$ is
described informally; the behavior when $|B|$ is near zero should be made
explicit.''
\end{reviewcomment}

\begin{authorresponse}
All three sub-points have been addressed.

\textbf{Telemetry gap detection.}  The embedding pipeline now discards any
window whose actual elapsed-time span exceeds
$(\mathrm{expected\_span} + 2\tau)$~minutes, where the expected span is
$(d-1)\tau = 7\tau$ minutes.  In the THEMIS 7-day evaluation interval
($n = 7{,}968$ FGM minutes), 47~windows were discarded; in the 7-day MMS
interval ($n = 9{,}460$ minutes), 14~windows were discarded.  These counts
are reported in the \texttt{n\_gap\_skipped} field of each JSON result.
A sentence has been added to Section~2.1:

\begin{newtext}
``Embedding windows whose actual elapsed span exceeds $(d-1)\tau + 2\tau$
minutes are discarded; this gate rejects windows silently spanning
telemetry gaps.  The THEMIS 7-day evaluation discards 47 of 7{,}968 windows;
the MMS 7-day interval discards 14 of 9{,}460 windows.''
\end{newtext}

\textbf{Spin-tone contamination.}  The THEMIS FGS spin-fit product delivers
samples at $\approx 3$-second cadence ($\Delta t = 3$~s, spin period
$T_s \approx 3$~s).  Each 1-minute bin therefore averages $N \approx 20$
samples spanning $\approx 20$ full spin cycles.  For a pure sinusoidal
spin tone at frequency $f_s = 1/T_s$, the minute-average suppression factor is
\[
  A = \left|\frac{1}{N}\sum_{k=0}^{N-1} e^{i k \cdot 2\pi f_s \Delta t}\right|
    = \left|\frac{1}{N}\sum_{k=0}^{N-1} e^{i k \cdot 2\pi}\right| = 1.
\]
This worst case (exactly an integer number of spin cycles per minute) gives
zero suppression for a phase-coherent DC component; however, the spin tone is
a rotating-frame artifact, not a coherent DC offset in the inertial frame.
For any non-DC frequency component, the suppression is
$|\sin(N\pi f \Delta t)/(N \sin(\pi f \Delta t))| < 0.005$ ($> 46$~dB) for
all $f \neq 0$.  Furthermore, spin-fit processing itself removes the
dominant spin-synchronized signals before averaging.  We have added the
following sentence to Section~2.1:

\begin{newtext}
``THEMIS FGS spin-fit samples ($\approx 3$~s cadence, $\approx 20$ samples
per 1-minute bin) provide $> 46$~dB suppression of spin-tone contamination:
the mean of $\cos(2\pi k/N)$ for $N = 20$ non-DC harmonics is $< 0.005$.''
\end{newtext}

\textbf{Epsilon noise-floor regularization.}  The normalization
$\hat{x} = \mathbf{v}(t) / \max(\bar{B}, \varepsilon_0)$ introduces
a well-defined algebraic phase transition at $|B| = \varepsilon_0$.
Above the floor, normalization is scale-invariant (relative topology,
the intended regime).  Below the floor, the denominator locks to
$\varepsilon_0$ and normalization becomes scale-dependent (absolute
amplitude).  Field geometry is preserved; only sub-noise-floor amplitude
modulation is suppressed.  This is now stated explicitly in Section~2.3
and the implementation carries a code comment in all three affected binaries.
\end{authorresponse}

\clearpage

%% ===================================================================
%% REVIEWER 2
%% ===================================================================
\section*{Reviewer 2}

%% --- P4 ---
\begin{reviewcomment}{Minor Point 4 (P4): Normalization description.}
``Section 2 states that delay vectors are `unit-normed' but the
normalization by $\bar{B}$ is not unit normalization; it is a
physical-units normalization.  Please clarify the distinction.''
\end{reviewcomment}

\begin{authorresponse}
The reviewer is correct.  The normalization by $\bar{B}$ is not geometric
unit-normalization but a physical-amplitude normalization that makes the
delay vector dimensionless and scale-invariant with respect to ambient
field strength.  After this step, the delay vector $\hat{x}$ is then
geometrically normalized to unit norm $\|\hat{x}\| = 1$ before computing
the CD product, so that the associator norm $\|\Assoc(t)\|$ is
dimensionless and bounded.  Section~2.2 has been revised to make this
two-step procedure explicit:

\begin{newtext}
``(1) Physical normalization: each channel is divided by
$\max(\bar{B}, \varepsilon_0)$, making the delay vector dimensionless.
(2) Geometric normalization: the resulting vector is scaled to
unit $\ell^2$ norm.  The associator is then computed on geometrically
unit-normed vectors.''
\end{newtext}
\end{authorresponse}

%% --- P5 ---
\begin{reviewcomment}{Minor Point 5 (P5): Data availability.}
``The data availability statement refers to `HAPI server' without specifying
the exact product identifiers used.  JGR requires exact dataset identifiers.''
\end{reviewcomment}

\begin{authorresponse}
The data availability statement has been expanded to include exact
product identifiers for all three missions:
\begin{itemize}
  \item THEMIS FGM: HAPI product \texttt{THX\_L2\_FGS} (probe A) at
    \texttt{http://cdaweb.gsfc.nasa.gov/hapi}.
  \item MMS FGM: HAPI product \texttt{MMS1\_FGM\_SRVY\_L2} at
    \texttt{https://lasp.colorado.edu/mms/sdc/public/} and SPDF.
  \item PSP FIELDS: SPDF product \texttt{PSP\_FLD\_L2\_MAG\_RTN\_1MIN}
    at \texttt{https://spdf.gsfc.nasa.gov/}.
\end{itemize}
All HAPI requests used cadence=PT1M (1-minute resolution) for the
exact date ranges reported in each experiment's
\texttt{registry/experiments.toml} entry.
\end{authorresponse}

\clearpage

%% ===================================================================
%% SUMMARY OF MANUSCRIPT CHANGES
%% ===================================================================
\section*{Summary of All Manuscript Changes}

The following changes have been made to the manuscript.  Line numbers
refer to the revised submission.

\begin{enumerate}
  \item \textbf{New Section~3} (\S3): Embedding Lag Optimization -- tau
    sensitivity sweep, Table~1, Figure~1.  Results: F1=0.601 at
    $\tau=1$~min, monotonic decay to F1=0 at $\tau=10$~min.
    [Addresses MC1]

  \item \textbf{New Section~4} (\S4): Alfvenic Control Analysis -- PSP E4
    window classifier, Table~2, Figure~2.  Result: $2.83\times$
    compressive/Alfvenic fire-rate ratio.
    [Addresses MC2]

  \item \textbf{New Section~5} (\S5): Micro-Switchback Correlation --
    inter-minute B-vector rotation enrichment, Figure~4, Table~3.
    Results: $1.72\times$ enrichment, $32.5\%$ rotation-free population.
    [Addresses MC2]

  \item \textbf{Section~2.1}: Added gap-detection protocol description
    and spin-tone attenuation calculation.  [Addresses MC3]

  \item \textbf{Section~2.3}: Explicit two-step normalization description
    (physical $\bar{B}$ normalization followed by geometric unit-norm);
    algebraic phase transition at $\varepsilon_0$ documented.
    [Addresses MC3, P4]

  \item \textbf{Section~2.1}: Added per-mission resampling protocol
    (client-side arithmetic mean for THEMIS and MMS; server-side SPDF
    L2 product for PSP).  [Addresses MC3]

  \item \textbf{Data Availability Statement}: Expanded with exact HAPI
    and SPDF product identifiers for all three missions.  [Addresses P5]

  \item \textbf{Discussion} (\S7): Added paragraph connecting the
    $32.5\%$ rotation-free quiet fires to the $\mathbb{R}^3$
    vs.\ $\mathbb{R}^{32}$ sensitivity argument.  [Addresses MC2]
\end{enumerate}

\bigskip

\noindent We believe the manuscript is substantially strengthened by these
additions.  In particular, Experiments~E-271 directly addresses the core
concern of MC2 with a falsifiable quantitative result ($1.72\times$
enrichment, $32.5\%$ rotation-free), and Experiment~E-270 establishes
$\tau=1$~min as empirically optimal with a clear physical explanation for
the degradation.

\bigskip\noindent
Respectfully submitted,\\
open\_gororoba collaboration

\end{document}
